\documentclass[10pt,a4paper]{amsart}
\usepackage[margin=19mm]{geometry}
\usepackage{amsmath,amssymb,mathtools}
\usepackage[scr=rsfs,cal=euler]{mathalfa}
\usepackage{xfrac}
\usepackage[unicode,hidelinks,bookmarks=false]{hyperref}
\newcommand{\ie}{i.e.\ }
\newcommand{\cf}{cf.\ }
\newcommand{\resp}{respectively\ }
\newcommand{\Subsection}{\subsection}
\providecommand{\nobreakdash}{\nobreak\hbox{-}}
\setlength{\emergencystretch}{2em}
\multlinegap0pt
\usepackage{bm}
\usepackage{bbm}
\usepackage{dsfont}
\usepackage{xspace}
\usepackage{cancel}
\usepackage{mathtools}
\usepackage{amsthm}
\makeatletter
\@ifundefined{thm}{\newtheorem{thm}{Thm}}{}
\@ifundefined{theorem}{\newtheorem{theorem}[thm]{Theorem}}{}
\makeatother
\newcommand{\bfA}[1]{{\textbf A}(#1)}
\title[Numerical Optimization for Tensor Disentanglement]{Numerical Optimization for Tensor Disentanglement}
\author{Julia Wei, Alec Dektor, Chungen Shen, Zaiwen Wen, Chao Yang}
\date{Source: arXiv:2508.19409v1 (CC BY 4.0)}
\begin{document}
\maketitle

\noindent\emph{Source and licence.} Numerical Optimization for Tensor Disentanglement. Version arXiv:2508.19409v1 (CC BY 4.0), distributed under the Creative Commons Attribution 4.0 International licence, as recorded in the arXiv licence metadata for this exact version and shown on the abstract page https://arxiv.org/abs/2508.19409v1. Full rights notice: licenses/arxiv-2508.19409-CC-BY-4.0.txt.\par\smallskip

\noindent\textit{Editorial scope.} Counted unit: the Theorem whose source label is \texttt{thm:euclidean\_grad} in the article (source0.tex lines 270--276, source label \texttt{thm:euclidean\_grad}), delivered together with the article's own complete original proof of it (source0.tex lines 278--310, one \texttt{proof} environment of 33 lines). No mathematics is written, repaired or re-derived: statement and proof are the article's own text line for line. Required context retained uncounted: eq:g\_def (lines 188--190). Every definition, display and label of those lines is retained unchanged. Omitted from this package and not redistributed: 1--187, 191--269, 277--277, 311--927. Nothing inside a retained excerpt is omitted or altered: every retained body is delivered exactly as the article's lines are written, so no omission receipt of any kind accompanies this package. Citations: the retained excerpts carry no citation command at all, so no bibliography entry of the article is transcribed and no reference is redistributed. Reference adaptation: none was needed and none was made; every reference inside the delivered unit resolves inside it, so no unresolved reference or citation is produced. Printed slips kept literal: no printed word, symbol, hypothesis or number of the counted statement or of its proof was altered. Layout and class: the article's own class is \texttt{sn-jnl} with packages including lscape, amssymb, palatino, geometry, longtable, amsthm, amsfonts, amsmath, latexsym, epsfig, graphicx, xcolor, epstopdf, subfigure; the wrapper is the lane's pinned \texttt{amsart} editorial wrapper at 10pt on A4, which restates the theorem-like environments the retained text needs and adds the article's own macro definitions that the retained text needs (\texttt{\textbackslash bfA}), with extra wrapper packages (bm, bbm, dsfont, xspace, cancel, mathtools). The delivered numbering is the unit's own. Assets: no retained span contains a figure environment, a graphics-inclusion command, a PDF-page inclusion, a table or a TikZ picture; no image file is copied into this package, the package declares no asset dependency and no image file is redistributed with it. Licence: version arXiv:2508.19409v1 of this article is distributed under the Creative Commons Attribution 4.0 International licence, as recorded in the arXiv licence metadata for this exact version and shown on the abstract page https://arxiv.org/abs/2508.19409v1. A full rights notice is delivered at licenses/arxiv-2508.19409-CC-BY-4.0.txt. Characters: every retained excerpt is pure ASCII, so the delivered engine, XeLaTeX with its default Latin Modern text font, reports no missing character.
\begin{equation} \label{eq:g_def}
    f(Q) = \sum_{i=1}^m \phi\left(\sigma_i(\bfA{QX})\right),
\end{equation}

\begin{theorem} \label{thm:euclidean_grad} 
    If $\bfA{QX} = U(Q) \Sigma(Q) V(Q)^{\top}$ is the singular value decomposition, then the gradient of $f$ in \eqref{eq:g_def} with respect to $Q$ is 
    \begin{equation} \label{eq:euclidean_grad}
    \nabla_Q f(Q) = {\bf A}^{-1}\left(U\phi'(\Sigma) V^{\top} \right) X^{\top}, 
    \end{equation}
    where the derivative $\phi'$ of $\phi$ is applied entry-wise to $\Sigma$. 
\end{theorem}

\begin{proof}
   We begin by taking the directional derivative of \eqref{eq:g_def} in the direction $E \in \mathbb{R}^{lr \times lr}$ and applying the chain rule to obtain 
   \begin{equation} \label{eq:g_diff}
   \begin{aligned}
       Df(Q)[E] &= \sum_{i=1}^m \phi'\left(\Sigma_{ii}\right) D\Sigma_{ii}(Q)[E] \\
       &= \mbox{trace}\left(\phi'(\Sigma) D \Sigma(Q)[E]\right), 
   \end{aligned}
   \end{equation}
   where the second equality holds since $\phi'(\Sigma)$ is diagonal. 
   Next, we obtain an expression for the directional derivative $D\Sigma(Q)[E]$. To do so, we differentiate $\bfA{QX}$ 
   \begin{equation}
   \begin{aligned}
       D \bfA{QX}[E] &= \bfA{EX} \\
       &= \left(DU(Q)[E]\right) \Sigma V^{\top} + U \left(D\Sigma(Q)[E]\right) V^{\top} + U\Sigma \left(DV(Q)^{\top}[E]\right),
    \end{aligned}
   \end{equation}
   where we used the fact that ${\bf A}$ is linear to obtain the first equality and the SVD representation of $\bfA{QX}$ with the product rule to obtain the second equality. Multiplying the preceding equation on the left by $U^{\top}$ and on the right by $V$ we obtain 
   \begin{equation} \label{eq:sigma_diff}
    D \Sigma(Q)[E] = U^{\top} \bfA{EX} V,
   \end{equation}
   where we used the fact that $U$ and $V$ have orthonormal columns and that $U^{\top}\left(DU(Q)[E]\right) =0$ and $\left(DV(Q)^{\top}[E]\right)V=0$. 
   Substituting \eqref{eq:sigma_diff} into \eqref{eq:g_diff} we obtain 
   \begin{equation} \label{eq:cyclic_trace}
   \begin{aligned}
       Df(Q)[E] &= \mbox{trace}\left(\phi'(\Sigma) U^{\top} \bfA{EX}V\right) \\
       &= \mbox{trace}\left(V \phi'(\Sigma) U^{\top} \bfA{EX}\right) \\
       &= \left\langle U\phi'(\Sigma)V^{\top}, \bfA{EX} \right\rangle \\
       &= \left\langle {\bf A}^{-1}\left( U\phi'(\Sigma)V^{\top}\right), EX \right\rangle \\
       &= \left\langle {\bf A}^{-1}\left( U\phi'(\Sigma)V^{\top}\right)X^{\top}, E \right\rangle
    \end{aligned}
   \end{equation}
   where we used the matrix inner product defined as $\langle X,Y\rangle=\mbox{trace}(X^{\top}Y)$ for any matrices $X,Y$ of the same dimension. The gradient $\nabla_Q f(Q)$ is the unique matrix satisfying $Df(Q)[E] = \langle \nabla_Q f(Q),E\rangle$ for all matrices $E$ of appropriate dimension. Comparing this with the last line in \eqref{eq:cyclic_trace} completes the proof. 
\end{proof}

\end{document}
